\section*{Part VIII: A Founder's Operating Method}

\subsection*{109. Step 1: Define the Larger Market Without Choosing the Entry Yet}

Begin with the opportunity landscape created in Chapter 5. List the major exchange situations, not just product categories. Identify who buys, who uses, what jobs motivate payment, which alternatives dominate, and where attractive economics might exist. Preserve the broad map while narrowing operations.

For example, professional AI capability development could include individual self-study, employer learning budgets, internal training teams, independent technical mentoring, college-integrated electives, research seminars, consulting, and software for evaluation or assessment. These are not interchangeable ``AI education customers.''

\subsection*{110. Step 2: Produce 6 to 12 Candidate Customer Situations}

Create a list of plausible segments using real observable differences in buying behavior. Do not invent exotic combinations to appear original. Consider trigger, job, capability level, payer, buying setting, and constraints. Each candidate should be describable in a sentence a knowledgeable outsider could use to identify actual people.

Record what makes each segment genuinely distinct. If 2 candidates use the same buyer, budget, offer, channel, and delivery process, combine them. If 1 broad segment includes self-pay learners and institutions with purchasing committees, split it.

\subsection*{111. Step 3: Screen for Fatal Constraints}

Run the hard gates: legitimacy, plausibility of demand, reachability, deliverability, survivability, and testability. Any proposal that requires unapproved use of private data, unsupported accreditation claims, irresponsible debt-funded enrollment, or a financial exposure the founder cannot bear must be redesigned or deferred.

Record failed candidates rather than deleting them. A future team or partnership may change the constraints, but it does not change the present facts.

\subsection*{112. Step 4: Investigate 3 Serious Finalists in Parallel, Then Choose 1}

For the top 3 plausible groups, conduct real customer research. Use Chapter 6's methods: recent decision histories, actual alternatives, attempted purchases, workarounds, and buyer-process reconstruction. Seek contradictory evidence, not just supporters. Observe available artifacts where permitted and distinguish reports from verified behavior.

A founder with limited time should not endlessly interview every segment. The practical stopping question is: \emph{Have we learned enough to specify a paid experiment whose result could change our choice?} The required number of interviews depends on heterogeneity and stakes; there is no magic saturation count.

\subsection*{113. Step 5: Build a Beachhead Dossier}

For the preferred group, prepare 1 concise record including the target and anti-target, urgent job, buying trigger, default alternative, payer, budget path, acquisition setting, minimum proof, initial product, delivery constraints, unit economics, risks, and future adjacent groups. Every important assertion gets a confidence label and a plan to verify it.

If the dossier contains nothing but assumptions, that is not failure; it reveals the tasks to perform before investing in an expensive launch.

\subsection*{114. Step 6: Design the Smallest Complete Transaction}

A complete transaction includes a buyer who can pay, a truthful defined offer, delivery of actual value, a fair service or refund policy, and an outcome that can be inspected. Do not pre-sell beyond safe capacity or use vague promises to secure deposits.

The initial product should reduce a consequential uncertainty without making a permanent, high-capital commitment. Where feasible, price it so that payment becomes evidence of demand while the buyer receives a fair standalone benefit.

\subsection*{115. Step 7: Create a 1-Page Economic Model}

Estimate realized selling price, discounts, payment delays, refunds, variable delivery expenses, cohort or project fixed costs, customer acquisition cash, founder hours, and the capacity ceiling. Prepare pessimistic, base, and optimistic scenarios. Put the funding requirement on the model: a business can be profitable in an income statement and still need cash before delivery.

Use normalized labor costs to test long-run viability while separately tracking actual cash burn. If the initial economics work only because the founder supplies unlimited free labor, name that limitation explicitly.

\subsection*{116. Step 8: Select the Experimental Commitment Level}

Different decisions demand different evidence. A 2-hour paid workshop might be a low-risk experiment. A 12-month enterprise service contract may require security and operational controls before the first sale. A new classroom lease or residential campus investment demands much stronger demand and financial proof.

State the exact amount of money, founder time, staff obligation, and reputation risk you are willing to expose \emph{before} the test begins. This turns a tempting idea into an affordable learning bet.

\subsection*{117. Step 9: Measure the Whole Loop}

Record the journey from reach to actual outcome: eligible prospects approached, substantive conversations, qualified buyers, real offers, paid decisions, delivery success, support time, refunds, independently checkable results, and honest referrals. Include acquisition source and relationship dependence.

For a course, track completion and assessed work, not simply attendance. For an enterprise pilot, track acceptance of the actual deliverables, not just invoice size. For a marketplace, track completed transactions and successful matches, not sign-ups. These behaviors establish different kinds of evidence.

\subsection*{118. Step 10: Compare Evidence Against Predeclared Alternatives}

Before seeing results, specify conditions for repeat, redesign, change segment, or stop. For example, continue if paid demand exists and delivery is high quality at viable contribution; redesign if eligible prospects see value but reject the time commitment; change payer if engineers want training but consistently require employer funding; stop if there is neither urgent need nor a feasible transaction path.

Do not revise the success definition merely to protect the original idea. Small samples create uncertainty, so decisions may be conditional rather than binary. A useful record shows what the result actually established, what it did not, and the next decisive experiment.

\subsection*{119. Step 11: Concentrate Until a Transferable Asset Exists}

After a promising first experiment, repeat the same buyer/job combination where feasible. Improve acquisition, onboarding, product clarity, delivery reliability, and cost structure. Turn founder knowledge into process and evidence another team member can inspect. Collect references with consent and contextual accuracy.

The objective is not to refuse change. It is to extract enough genuine learning from a coherent customer group that the next decision is easier and more informed than the first.

\subsection*{120. Step 12: Investigate the Next Market Before Launching It}

Create an adjacency hypothesis. Identify the real transferred assets and the new unknowns. If the potential second market has a different payer, budget cycle, security regime, or faculty requirement, investigate those explicitly. Design a small experiment that can be funded without degrading the core.

A company that launches another product solely because it has spare creative energy has not demonstrated a scaling strategy. A company that can show why the next customer is cheaper to reach, easier to trust, or more reliably served because of the first may have one.

\subsection*{121. The Monthly Beachhead Review}

Once a month, write answers to 7 questions:

\begin{enumerate}
    \item Who exactly bought, and which triggering event caused the decision?
    \item Who did not buy, and what did they do instead?
    \item Which acquisition path reached eligible buyers without excessive founder time?
    \item What actually improved for customers, and what evidence supports that inference?
    \item What was the contribution and cash flow after realistic operating costs?
    \item Which repeatable capability, proof asset, or relationship now exists that did not exist last month?
    \item Which assumption has become less credible, and what decision should change as a result?
\end{enumerate}

These questions keep the company honest about the difference between progress, activity, and speculation.

\subsection*{122. The Conclusion: Win the Right to Be Broader}

A founder does not become strategically focused by uttering the phrase ``we serve a niche.'' Focus is the real allocation of scarce resources toward a specific group of customers, a specific costly job, and a specific mechanism of value creation. It is an operating choice with consequences: some customers must be refused, some products postponed, and some attractive opportunities deliberately left unpursued.

The benefit is that the first market can become knowable. You can watch purchases, understand objections, observe customer work, measure delivery, and correct the product. If your promise works repeatedly, the market begins to generate assets: cash, reputation, processes, talent, knowledge, and perhaps distribution. Those assets can make the next group easier to serve. If they do not, the first market may still be a good independent business, but it has not yet furnished a compounding expansion path.

The hardest decision is often not selecting which market sounds largest. It is selecting the \textbf{smallest strategically consequential problem that a new organization can actually solve and sell today.}

\begin{quote}
\emph{Do not begin by being vaguely useful to everyone. Begin by being unusually useful to someone, prove that the economics and the outcome survive repetition, and expand only along paths where what you have learned makes the next success more likely.}
\end{quote}

That is how an ambitious founder preserves the large dream without allowing it to consume the first small business before the business has had a chance to exist.
